% Part: sets-functions-relations
% Chapter: functions
% Section: inverses

\documentclass[../../../include/open-logic-section]{subfiles}

\begin{document}

\olfileid{sfr}{fun}{inv}
\olsection{વિધેયોનાં પ્રતિવિધેયો}

\begin{explain}
આપણે વિધેયોને નિરૂપણો તરીકે વિચારીએ છીએ. તેથી વિધેયો
વિશે એક સ્વાભાવિક પ્રશ્ન એ છે કે નિરૂપણને “ઊલટાવી”
શકાય કે નહીં. દાખલા તરીકે, અનુગામી વિધેય
$f(x) = x + 1$ને આ અર્થમાં ઊલટાવી શકાય કે વિધેય
$g(y) = y - 1$ એ $f$ની ક્રિયાને “પાછી ફેરવે” છે.

પરંતુ સાવચેત રહેવું જોઈએ. $g$ની વ્યાખ્યા વિધેય
$\Int \to \Int$ વ્યાખ્યાયિત કરે છે, છતાં તે
\emph{વિધેય} $\Nat
\to \Nat$ વ્યાખ્યાયિત કરતી નથી, કારણ કે
$g(0) \notin \Nat$. એટલે સરળ કિસ્સામાં પણ વિધેયને
ઊલટાવી શકાય કે નહીં તે તરત સ્પષ્ટ થતું નથી;
તે પ્રદેશ અને સહપ્રદેશ પર આધાર રાખી શકે છે.

વિધેયના પ્રતિવિધેયનો ખ્યાલ આ વાતને વધુ ચોક્કસ બનાવે છે.
\end{explain}

\begin{defn}
વિધેય $g \colon B \to A$ એ વિધેય $f
\colon A \to B$નું \emph{પ્રતિવિધેય} છે જો બધા
$x \in A$ અને $y \in B$ માટે
$f(g(y)) = y$ અને $g(f(x)) = x$ થાય.
\end{defn}

જો $f$નું પ્રતિવિધેય~$g$ હોય, તો ઘણી વાર $g$ને
બદલે $f^{-1}$ લખીએ છીએ.

\begin{explain}
હવે વિધેયોને ક્યારે પ્રતિવિધેયો હોય છે તે નક્કી કરીશું.
$f\colon A \to B$ના પ્રતિવિધેય માટેનો યોગ્ય ઉમેદવાર
નીચે મુજબ “વ્યાખ્યાયિત” $g\colon B \to A$ છે:
\[
g(y) = \text{$f(x) = y$ થાય એવો “તે એક જ” $x$.}
\]
પરંતુ “વ્યાખ્યાયિત” (અને “તે એક જ”)ની આસપાસનાં
અવતરણચિહ્નો સૂચવે છે કે આ વ્યાખ્યા નથી. ઓછામાં ઓછું,
તે સંપૂર્ણ સામાન્યતાથી હંમેશાં કામ નહીં કરે. કારણ કે
આ વ્યાખ્યા વિધેય નક્કી કરે તે માટે $f(x) =
y$ થાય એવો એક અને માત્ર એક~$x$ હોવો જોઈએ---$g$નું
નિર્ગત અનન્ય રીતે નક્કી થવું જોઈએ. વધુમાં, દરેક
$y \in B$ માટે તે નક્કી થવું જોઈએ. જો $x_1$ અને
$x_2 \in
A$ એવા હોય કે $x_1 \neq x_2$ છતાં
$f(x_1) = f(x_2)$ હોય, તો $y = f(x_1) = f(x_2)$
માટે $g(y)$ અનન્ય રીતે નક્કી નહીં થાય. અને જો
$f(x) = y$ થાય એવો કોઈ~$x$ જ ન હોય, તો
$g(y)$ બિલકુલ નક્કી થતું નથી. બીજા શબ્દોમાં,
$g$ વ્યાખ્યાયિત થાય તે માટે $f$ એ !!{injective}
અને !!{surjective} બંને હોવું જોઈએ.

ધીમે ધીમે આગળ વધીએ. પ્રશ્નને બે ભાગમાં વહેંચીએ:
વિધેય~$f\colon A \to B$ આપેલું હોય ત્યારે
$g(f(x)) = x$ થાય એવું વિધેય $g\colon B \to A$
ક્યારે હોય? આવું $g$ એ $f$ની ક્રિયાને “પાછી ફેરવે” છે;
તેને $f$નો \emph{ડાબી બાજુનો વ્યસ્ત} કહે છે.
બીજું, $f(h(y)) = y$ થાય એવું વિધેય
$h\colon B \to A$ ક્યારે હોય? આવા $h$ને
$f$નો \emph{જમણી બાજુનો વ્યસ્ત} કહે છે---$f$
એ $h$ની ક્રિયાને “પાછી ફેરવે” છે.
\end{explain}

\begin{prop}
જો $f\colon A \to B$ એ !!{injective} હોય અને તેનો પ્રદેશ અખાલી હોય, તો
$f$નો \emph{ડાબી બાજુનો વ્યસ્ત}~$g\colon B \to A$
હોય છે, જેથી બધા $x
\in A$ માટે $g(f(x)) = x$ થાય.
\end{prop}

\begin{proof}
ધારો કે $f\colon A \to B$ એ !!{injective} છે અને તેનો પ્રદેશ અખાલી છે.
કોઈ $y \in B$ લો. જો $y \in \ran{f}$ હોય,
તો $f(x) = y$ થાય એવો $x \in A$ છે.
$f$ એ !!{injective} હોવાથી આવો માત્ર એક~$x \in A$
છે. તેથી $g(y) = x$ વ્યાખ્યાયિત કરી શકીએ, એટલે કે
$g(y)$ એ $f(x)
= y$ થાય એવો “તે એક જ” $x \in A$ છે.
જો $y \notin \ran{f}$ હોય, તો તેને કોઈ પણ~$a \in A$
પર મોકલી શકીએ. એટલે એક $a \in A$ પસંદ કરી
$g\colon B \to A$ને આમ વ્યાખ્યાયિત કરી શકીએ:
\[
g(y) = \begin{cases}
    x & \text{જો $f(x) = y$ હોય}\\
    a & \text{જો $y \notin \ran{f}$ હોય.}
\end{cases}
\]
તે બધા $y \in B$ માટે વ્યાખ્યાયિત છે, કારણ કે
આવા દરેક $y \in \ran{f}$ માટે $f(x) = y$ થાય
એવો ચોક્કસ એક $x \in A$ છે. વ્યાખ્યા પ્રમાણે,
જો $y = f(x)$ હોય, તો $g(y) = x$ થાય,
એટલે કે $g(f(x)) = x$.
\end{proof}
\sourcecorrection{OLFUN-001}{મૂળની \texttt{inverses.tex}, પંક્તિઓ 62--84માં
પ્રદેશ અખાલી હોવાની શરત વિના આ વિધાન ખોટું છે: ખાલી ગણથી એક-ઘટકવાળા
ગણમાંનું ખાલી વિધેય એક-એક છે, પરંતુ તેનો ડાબી બાજુનો વ્યસ્ત નથી.
સરળ સુધારેલા વિધાનમાં પ્રદેશ અખાલી હોવાની શરત ઉમેરી છે. ચોક્કસ સામાન્ય
શરત એ છે કે પ્રદેશ અખાલી હોય અથવા સહપ્રદેશ ખાલી હોય.}

\begin{prob}
બતાવો કે જો $f\colon A \to B$નો ડાબી બાજુનો
વ્યસ્ત~$g$ હોય, તો $f$ એ !!{injective} છે.
\end{prob}

\begin{prop}
જો $f\colon A \to B$ એ !!{surjective} હોય, તો
$f$નો \emph{જમણી બાજુનો વ્યસ્ત}~$h\colon B \to A$
હોય છે, જેથી બધા~$y \in B$ માટે $f(h(y)) =
y$ થાય.
\end{prop}

\begin{proof}
ધારો કે $f\colon A \to B$ એ !!{surjective} છે.
કોઈ $y \in
B$ લો. $f$ એ !!{surjective} હોવાથી
$f(x_y)
= y$ થાય એવો $x_y \in A$ છે.
તેથી $h(y) = x_y$ વ્યાખ્યાયિત કરી શકીએ, એટલે કે
દરેક $y \in B$ માટે $f(x) = y$ થાય એવો કોઈ
$x \in A$ પસંદ કરીએ; $f$ એ !!{surjective}
હોવાથી પસંદ કરવા માટે હંમેશાં ઓછામાં ઓછો એક હોય છે.
\footnote{$f$ એ !!{surjective} હોવાથી દરેક~$y \in B$
માટે ગણ $\Setabs{x}{f(x) = y}$ ખાલી નથી.
$h$ની આપણી વ્યાખ્યા માટે આ દરેક ગણમાંથી એક જ $x$
પસંદ કરવો પડે છે. આ હંમેશાં શક્ય છે તે હકીકતમાં
સ્પષ્ટ નથી---આવી પસંદગીઓ કરવાની શક્યતા એક
સ્વયંસિદ્ધિ તરીકે માની લેવામાં આવે છે.
બીજા શબ્દોમાં, આ વિધાન કહેવાતી પસંદગીની સ્વયંસિદ્ધિ
ધારે છે. આ મુદ્દાને આપણે
\oliflabeldef{sth:choice::chap}{\olref[sth][choice][]{chap}માં ફરી જોઈશું}{અહીં વિગત વિના જવા દઈશું}.
જોકે, ઘણા ચોક્કસ કિસ્સામાં, દાખલા તરીકે $A = \Nat$
હોય કે સાન્ત હોય, અથવા $f$ એ !!{bijective} હોય,
ત્યારે પસંદગીની સ્વયંસિદ્ધિ જરૂરી નથી. (ખાસ કરીને
$f$ એ !!{bijective} હોય ત્યારે દરેક $y \in B$
માટે ગણ $\Setabs{x}{f(x) = y}$માં ચોક્કસ એક
!!{element} હોય છે, એટલે પસંદગી કરવાની રહેતી જ નથી.)}
વ્યાખ્યા પ્રમાણે, જો $x = h(y)$ હોય,
તો $f(x) = y$ થાય; એટલે કે કોઈ પણ $y \in B$
માટે $f(h(y)) = y$.
\end{proof}

\begin{prob}
બતાવો કે જો $f\colon A \to B$નો જમણી બાજુનો
વ્યસ્ત~$h$ હોય, તો $f$ એ !!{surjective} છે.
\end{prob}

\begin{explain}
અગાઉની સાબિતીના વિચારોને ભેગા કરતાં હવે મળે છે કે
દરેક !!{bijection}ને પ્રતિવિધેય હોય છે, એટલે કે
એક જ વિધેય $f$નો ડાબી અને જમણી બંને બાજુનો વ્યસ્ત હોય છે.
\end{explain}

\begin{prop}\ollabel{prop:bijection-inverse}
જો $f\colon A \to B$ એ !!{bijective} હોય, તો એવું
વિધેય~$f^{-1}\colon B \to A$ છે કે બધા $x \in A$
માટે $f^{-1}(f(x)) = x$ અને બધા $y \in B$
માટે $f(f^{-1}(y)) = y$ થાય.
\end{prop}

\begin{proof}
સ્વાધ્યાય.
\end{proof}

\begin{prob}
\olref[sfr][fun][inv]{prop:bijection-inverse} સાબિત કરો.
તમારે $f^{-1}$ વ્યાખ્યાયિત કરવું, તે વિધેય છે તે
બતાવવું, અને તે $f$નું પ્રતિવિધેય છે તે બતાવવું પડે;
એટલે કે બધા $x \in A$ અને $y \in B$ માટે
$f^{-1}(f(x)) = x$ અને $f(f^{-1}(y)) = y$ થાય.
\end{prob}

\begin{explain}
પ્રતિવિધેયો મેળવવાની થોડી વધુ સામાન્ય રીત છે.
આપણે \olref[kin]{sec}માં જોયું કે દરેક વિધેય $f$
બધા $x \in A$ માટે $f'(x) = f(x)$ લઈને
!!a{surjection} $f'
\colon A \to \ran{f}$ ઉત્પન્ન કરે છે.
સ્પષ્ટ છે કે $f$ એ !!{injective} હોય, તો
$f'$ એ !!{bijective} છે; તેથી
\olref{prop:bijection-inverse} પ્રમાણે તેને અનન્ય
પ્રતિવિધેય હોય છે. સંકેતપ્રયોગમાં બહુ નાની છૂટ લઈને
આપણે ક્યારેક $f'$ના પ્રતિવિધેયને માત્ર
“$f$નું પ્રતિવિધેય” કહીએ છીએ.
\end{explain}

\begin{prop}\ollabel{prop:left-right}%
બતાવો કે જો $f\colon A \to B$નો ડાબી બાજુનો
વ્યસ્ત~$g$ અને જમણી બાજુનો વ્યસ્ત~$h$ હોય,
તો $h = g$ થાય.
\end{prop}

\begin{proof}
સ્વાધ્યાય.
\end{proof}

\begin{prob}
\olref[sfr][fun][inv]{prop:left-right} સાબિત કરો.
\end{prob}

\begin{prop}\ollabel{prop:inverse-unique}
દરેક વિધેય~$f$ને વધુમાં વધુ એક પ્રતિવિધેય હોય છે.
\end{prop}

\begin{proof}
ધારો કે $g$ અને $h$ બંને $f$નાં પ્રતિવિધેયો છે.
તો ખાસ કરીને $g$ એ $f$નો ડાબી બાજુનો વ્યસ્ત
અને $h$ જમણી બાજુનો વ્યસ્ત છે.
\olref{prop:left-right} પ્રમાણે $g = h$.
\end{proof}

\end{document}
